\documentclass[11pt]{amsart}
\usepackage[margin=1in]{geometry}
\usepackage{amssymb,amsmath,booktabs}
\usepackage[colorlinks=true,linkcolor=blue,citecolor=blue,urlcolor=blue,pdfusetitle]{hyperref}

\newtheorem{theorem}{Theorem}
\newtheorem{proposition}[theorem]{Proposition}
\newtheorem{lemma}[theorem]{Lemma}
\newtheorem{corollary}[theorem]{Corollary}
\theoremstyle{remark}
\newtheorem{remark}[theorem]{Remark}

\newcommand{\Z}{\mathbb{Z}}
\newcommand{\Lset}{\mathcal{L}}

\title[Exact two-adic counts for the hyperbolic sieve]{Dyadic Shells, Version 11.7.23:\\ Exact two-adic counts for Fermat-type quadratic sieves and the hyperbolic sieve}
\author{Arsen A. Khachatryan}
\address{SibarS Labs, SibarS Group of Companies}
\email{cto@sibars.net}
\thanks{ORCID 0009-0006-7629-3833. Preprint, Version 11.7.23, 8 November 2026, DOI \href{https://doi.org/10.5281/zenodo.24051028}{10.5281/zenodo.24051028}; not peer-reviewed.}
\date{2026-11-08}

\begin{document}

\begin{abstract}
For an odd integer $N$ let $C_m(N)$ be the number of residues $A \bmod 2^m$ for which $A^2-N$ is a square modulo $2^m$. These counts describe the two-adic stage of Fermat-type factorisation sieves. Through the identity $x^2-4kN=(ku-v)^2$ they also determine the size of Hittmeir's hyperbolic sieve sets $\Lset_{N,2^e,k}$, for which only computed values have been recorded. We determine $C_m(N)$ for every odd $N$ and every $m\ge1$: it depends only on $N\bmod 8$; it equals $2^{m-2}$ for $N\equiv3\pmod4$ and $2^{m-3}$ for $N\equiv5\pmod8$ (for $m\ge3$, resp.\ $m\ge5$); and for $N\equiv1\pmod8$ it equals four times the number of squares modulo $2^{m-4}$ ($m\ge4$), i.e.\ $(2^{m-3}+16)/3$ or $(2^{m-3}+20)/3$ for even $m\ge6$ or odd $m\ge5$. We prove $|\Lset_{N,2^e,k}|=C_e(kN)/2$ for odd $N,k$ and $e\ge1$; this gives proofs of the values for powers of two that Hittmeir recorded computationally, with a recursion exponent that differs from the one printed in the arXiv version of that paper. Consequently the two-adic pass proportion satisfies $C_m(N)/2^m>1/24$ for all odd $N$ and all $m$, and $1/24$ cannot be replaced by a larger constant. We also show that a sieve prime $\ell$ dividing the multiplier exactly once removes a proportion $1/\ell$ of the residues modulo $\ell^2$ and nothing more at higher powers. The reduction to $N\bmod8$ is formally verified in Lean~4 with Mathlib; the closed forms are not yet formally verified.
\end{abstract}

\maketitle

\section{Introduction}

Fermat's method, in Lawrence's form~\cite{lawrence}, searches for a decomposition $N=uv$ through the identity
\[
(au+bv)^2-(au-bv)^2=4abN ,
\]
testing candidates $x$ for which $x^2-4abN$ is a perfect square. Practical implementations first discard candidates for which $x^2-4abN$ is not a square modulo some small moduli~\cite{mckee}. Hittmeir~\cite{hittmeir2018,hittmeir2023} recast this filter as the \emph{hyperbolic sieve}: for a modulus $M$ coprime to $Nk$ put
\[
\Lset_{N,M,k}=\{\,kx+y \bmod M : xy\equiv N \pmod M\,\},
\]
so that $ku+v \bmod M\in\Lset_{N,M,k}$ for every co-divisor pair $(u,v)$ of $N$. For an odd prime $r$ he proved $|\Lset_{N,r,k}|=(r+(k^{-1}N\mid r))/2$~\cite[Lemma~3.4]{hittmeir2023}, showed that for odd $M$ the set coincides with $\{x : x^2-4kN \text{ is a square mod } M\}$~\cite[Lemma~4.2]{hittmeir2023}, determined $|\Lset_{N,r^e,k}|$ for odd prime powers~\cite[Theorem~3.5]{hittmeir2023}, and used these sets to refine the complexity of Fermat's method by a subexponential factor~\cite[Theorem~4.6]{hittmeir2023}. The same sets underlie the deterministic algorithms of Harvey and Hittmeir~\cite{harveyhittmeir} and the multi-target and quadratic-image refinements of Stănică, Mulder and Hittmeir~\cite{smh2026} and Stănică~\cite{stanica2026}, which use odd auxiliary primes. (Theorem and lemma numbers of~\cite{hittmeir2023} refer to arXiv:2205.10074v1; we could not access the typeset journal version.)

For powers of two, \cite[Remark~7.5]{hittmeir2023} records the sizes of $\Lset_{N,2^e,k}$ ``according to our computations'' and suggests that they could be proved by modifying the odd-prime argument. Independently, the author's preprint series on anchored Fermat-type screening~\cite{khach11517} proved, for odd $N\ge3$, only the $N$-dependent bound $C_m(N)/2^m\ge2^{-(2v+3)}$ with $v=v_2((N-1)/2)$, and left the limiting value $1/24$ for $N\equiv1\pmod8$ as a conjecture supported by computation.

\subsection*{Relation to Version 11.5.17}
This is Version~11.7.23 of the preprint series on dyadic shells in semiprime arithmetic (series concept DOI 10.5281/zenodo.20306132). It is a separate paper on the two-adic stage of the sieve and does not modify Version~11.5.17~\cite{khach11517} (17~October 2026, DOI 10.5281/zenodo.23432486), which remains the record of what was stated then: for odd $N\ge3$ the lower bound $\pi_m(N)\ge2^{-(2v+3)}$, $v=v_2((N-1)/2)$ (Proposition~4.2 there), and the value $1/24$ for $N\equiv1\pmod 8$ as a conjecture supported by computation. Theorem~\ref{thm:main} and Corollary~\ref{cor:density} below prove that value and replace the bound by an exact count. The sieve counts for odd primes and odd prime powers used in Version~11.5.17 (Propositions~4.1 and~4.3 there) are due to Hittmeir~\cite{hittmeir2023}, and the identity behind its Proposition~3.1 is Lawrence's~\cite{lawrence}.

\subsection*{What is new and what is not}
The criterion for squares modulo $2^m$ (Lemma~\ref{lem:squares}) and the count of squares modulo $2^n$ (Lemma~\ref{lem:sqcount}, \cite{stangl}) are classical. The hyperbolic sieve and its sizes for odd primes and odd prime powers are due to Hittmeir~\cite{hittmeir2023}. To our knowledge, after a search of the cited literature, the OEIS and general web search, the closed forms of Theorem~\ref{thm:main}, the identity of Theorem~\ref{thm:hyp} and the sharp constant $1/24$ have not been stated with proof before; an automated reading of Shparlinski's survey on modular hyperbolas~\cite{shparlinski} found no such count either. This is not an expert literature review, in particular of the literature on $2$-adic local densities, and we make no stronger novelty claim. Nothing here gives a new general-purpose factoring algorithm or affects properly generated RSA keys.

\subsection*{Results}
Let $N$ be odd and
\[
C_m(N)=\#\{A\in\Z/2^m\Z : A^2-N \text{ is a square modulo } 2^m\},\qquad \pi_m(N)=C_m(N)/2^m .
\]
Here and below $m$, $e$, $k_0$ denote exponents of $2$, and $M$ denotes a general modulus.

\begin{theorem}\label{thm:main}
Let $N$ be odd. Then $C_1(N)=C_2(N)=2$, and for $m\ge3$:
\begin{enumerate}
\item if $N\equiv3\pmod4$, then $C_m(N)=2^{m-2}$;
\item if $N\equiv5\pmod8$, then $C_3(N)=C_4(N)=4$ and $C_m(N)=2^{m-3}$ for $m\ge5$;
\item if $N\equiv1\pmod8$, then $C_3(N)=4$ and $C_m(N)=4\,s(m-4)$ for $m\ge4$, where $s(n)$ is the number of squares modulo $2^n$. Explicitly, $C_4(N)=4$ and
\[
C_m(N)=\frac{2^{m-3}+16}{3}\ (m\ge6\text{ even}),\qquad C_m(N)=\frac{2^{m-3}+20}{3}\ (m\ge5\text{ odd}).
\]
\end{enumerate}
\end{theorem}

\begin{corollary}\label{cor:density}
For every odd $N$ and every $m\ge1$, $\pi_m(N)>1/24$. Moreover $\lim_{m\to\infty}\pi_m(N)$ equals $1/4$, $1/8$ or $1/24$ according as $N\equiv3\pmod4$, $N\equiv5\pmod8$ or $N\equiv1\pmod8$; in particular $\inf_{N,m}\pi_m(N)=1/24$.
\end{corollary}

\begin{theorem}\label{thm:hyp}
Let $N,k$ be odd and $e\ge1$. Then $|\Lset_{N,2^e,k}|=C_e(kN)/2$. Consequently $|\Lset_{N,2,k}|=|\Lset_{N,4,k}|=1$, $|\Lset_{N,8,k}|=1$ if $N\equiv3k\pmod4$ and $2$ otherwise, $|\Lset_{N,16,k}|=2$, $|\Lset_{N,32,k}|=2$ if $N\equiv5k\pmod8$ and $4$ otherwise, and for $e\ge5$
\[
|\Lset_{N,2^{e+1},k}|=
\begin{cases}
2\,|\Lset_{N,2^e,k}| & \text{if } N\equiv3k\pmod 4 \text{ or } N\equiv5k\pmod8,\\
2\bigl(|\Lset_{N,2^e,k}|-2^{\,e \bmod 2}\bigr) & \text{if } N\equiv k\pmod8.
\end{cases}
\]
\end{theorem}

\begin{remark}\label{rem:exponent}
In arXiv:2205.10074v1, the last case of Remark~7.5 is printed with the exponent $(e+1)\bmod 2$, and so is the analogous term of Theorem~3.5 (odd prime powers, case $(k^{-1}N\mid r)=1$). The values in Example~4.7 of the same text ($|\Lset_{N,2^e,1}|=2,4,4,6,8,14,24$ for $e=4,\dots,10$ and $N\equiv1\pmod8$; $4,16,106$ modulo $7,7^2,7^3$) agree with the exponent $e\bmod2$, and so do all finite tests of Section~\ref{sec:numerics}; with $(e+1)\bmod2$ none of the tested steps agrees. For the two-adic case, Theorem~\ref{thm:hyp} proves the exponent $e\bmod 2$. We have not seen the typeset journal version~\cite{hittmeir2023} and found no erratum; we therefore attribute the printed exponent only to the arXiv text.
\end{remark}

\subsection*{Computational meaning}
Corollary~\ref{cor:density} concerns the two-adic filter alone: among residues $A\bmod 2^m$ distributed uniformly, the condition ``$A^2-N$ is a square modulo $2^m$'' keeps a proportion $\pi_m(N)>1/24\approx4.17\%$, whatever $m$ is. It says nothing about a sieve composed of several moduli, whose survivor proportion is the product of local proportions only under an independence assumption that fails for structured candidate sets; nor does a sharper local density by itself imply a lower total bit cost. For odd $N\ge3$ the bound $2^{-(2v+3)}$, $v=v_2((N-1)/2)$, proved in~\cite{khach11517} remains valid but can be far weaker (e.g.\ $2^{-15}$ for $v=6$); for $N=1$ that bound is not defined.

\section{Squares modulo powers of two}

\begin{lemma}\label{lem:squares}
Let $n\ge1$ and $x\in\Z$. Then $x$ is a square modulo $2^n$ if and only if either $2^n\mid x$, or $x=2^{2i}w$ with $w$ odd, $2i<n$ and $w\equiv1\pmod{2^{\min(3,\,n-2i)}}$.
\end{lemma}

\begin{proof}
For odd $w$ and $t\ge1$, $w$ is a square modulo $2^t$ if and only if $w\equiv1\pmod{2^{\min(3,t)}}$: necessity because odd squares are $\equiv1\pmod 8$, sufficiency by lifting square roots modulo powers of two (equivalently, an odd $w\equiv1\pmod8$ is a square in $\Z_2$~\cite[Ch.~II, Thm.~4]{serre}). Now let $x\not\equiv0\pmod{2^n}$ with $v_2(x)=j<n$. If $x\equiv y^2\pmod{2^n}$ then $v_2(y^2)=j$: otherwise $v_2(x-y^2)=\min(j,v_2(y^2))\le j<n$, contradicting $2^n\mid x-y^2$. Hence $j=2i$ is even and $y=2^iz$ with $z$ odd, and $z^2\equiv w\pmod{2^{n-2i}}$. Conversely, if $z^2\equiv w\pmod{2^{n-2i}}$ then $(2^iz)^2\equiv x\pmod{2^n}$.
\end{proof}

For $t\ge1$ let $f(t)$ be the number of odd residues modulo $2^t$ that are squares modulo $2^t$. By the first sentence of the proof, $f(1)=f(2)=1$ and $f(t)=2^{t-3}$ for $t\ge3$.

\begin{lemma}[{\cite{stangl}}]\label{lem:sqcount}
The number $s(n)$ of squares modulo $2^n$ satisfies $s(0)=1$ and
\[
s(n)=1+\sum_{0\le i<n/2} f(n-2i)=\frac{2^n+9-(-1)^n}{6}\qquad(n\ge1).
\]
\end{lemma}

\begin{proof}
By Lemma~\ref{lem:squares} the squares modulo $2^n$ are $0$ and the residues $2^{2i}w$ with $2i<n$ and $w$ an odd square modulo $2^{n-2i}$; for fixed $i$ there are $f(n-2i)$ of them. For the closed form: $s(1)=1+f(1)=2$, $s(2)=1+f(2)=2$, $s(3)=1+f(3)+f(1)=3$, $s(4)=1+f(4)+f(2)=4$; and for $n\ge5$, $s(n)-s(n-2)=f(n)=2^{n-3}$, which the closed form satisfies, since $(2^n-2^{n-2})/6=2^{n-3}$.
\end{proof}

\begin{lemma}[halving]\label{lem:halving}
For $n\ge2$ and $x\in\Z$, $4x$ is a square modulo $2^n$ if and only if $x$ is a square modulo $2^{n-2}$.
\end{lemma}

\begin{proof}
If $x\equiv z^2\pmod{2^{n-2}}$ then $4x\equiv(2z)^2\pmod{2^n}$. Conversely, if $4x\equiv y^2\pmod{2^n}$ then $4\mid y^2$, so $y=2z$ and $4x\equiv4z^2\pmod{2^n}$, i.e.\ $x\equiv z^2\pmod{2^{n-2}}$.
\end{proof}

\section{\texorpdfstring{Reduction to $N\bmod 8$}{Reduction to N mod 8}}

\begin{lemma}[unit-square invariance]\label{lem:scaling}
Let $R$ be a finite commutative ring, $N\in R$ and $u\in R^\times$. Then the number of $A\in R$ with $A^2-u^2N$ a square in $R$ equals the number of $A\in R$ with $A^2-N$ a square in $R$.
\end{lemma}

\begin{proof}
$A\mapsto u^{-1}A$ is a bijection of $R$, and $(u^{-1}A)^2-N=u^{-2}(A^2-u^2N)$; multiplication by the square of a unit preserves squareness in both directions.
\end{proof}

\begin{proposition}\label{prop:mod8}
If $N,N'$ are odd and $N\equiv N'\pmod 8$, then $C_m(N)=C_m(N')$ for every $m\ge1$.
\end{proposition}

\begin{proof}
$NN'\equiv N'^2\equiv1\pmod8$, so by Lemma~\ref{lem:squares} there is $y$ with $y^2\equiv NN'\pmod{2^m}$. Let $a$ be an inverse of $N'$ modulo $2^m$ and $u=ya$. Then $u^2N'\equiv NN'a^2N'\equiv N\pmod{2^m}$, and $u$ is a unit because $u^2N'\equiv N$ is. Apply Lemma~\ref{lem:scaling} with $R=\Z/2^m\Z$.
\end{proof}

Lemma~\ref{lem:scaling} and Proposition~\ref{prop:mod8} are formally verified (Section~\ref{sec:lean}). It therefore suffices to treat $N\in\{1,3,5,7\}$; negative $N$ are covered as well.

\section{\texorpdfstring{Proof of Theorem~\ref{thm:main}}{Proof of Theorem 1}}

\emph{The cases $m\le2$.} Every residue is a square modulo $2$, so $C_1=2$. Modulo $4$ the squares are $0,1$; if $N\equiv1\pmod4$ exactly the odd $A$ pass ($A^2-N\equiv0$), and if $N\equiv3\pmod4$ exactly the even $A$ pass ($A^2-N\equiv1$). Hence $C_2=2$. From now on $m\ge3$; oddness of $N$ is used throughout, since it makes $A^2-N$ odd for even $A$ and even for odd $A$.

\emph{Case $N\equiv3\pmod4$.} For odd $A$, $A^2-N\equiv 1-N\equiv2\pmod4$ has valuation $1$ and is not a square modulo $2^m$ (Lemma~\ref{lem:squares}). For $A=2a$, $A^2-N=4a^2-N$ is odd, hence a square modulo $2^m$ iff $4a^2\equiv N+1\pmod 8$. For $N\equiv3\pmod 8$ this means $a$ odd, for $N\equiv7\pmod8$ it means $a$ even. As $a$ runs over $\Z/2^{m-1}\Z$, half of the values qualify: $C_m(N)=2^{m-2}$.

\emph{Case $N\equiv5\pmod8$.} Take $N=5$. Even $A$ give $A^2-5\equiv3\pmod4$, never a square. For $A=2t+1$ we have $A^2-5=4T$ with $T=t(t+1)-1$ odd. By Lemmas~\ref{lem:halving} and~\ref{lem:squares}, $A$ passes iff $T\equiv1\pmod{2^{\min(3,m-2)}}$, i.e.\ $t(t+1)\equiv2\pmod{2^{\min(3,m-2)}}$. As $A$ runs over the odd residues modulo $2^m$, $t$ runs over $\Z/2^{m-1}\Z$. For $m=3$ the condition is vacuous: $C_3=4$. For $m=4$ it reads $t\equiv1,2\pmod4$: $C_4=4$. For $m\ge5$ it reads $t\equiv1,6\pmod8$, a proportion $1/4$ of $\Z/2^{m-1}\Z$: $C_m=2^{m-3}$.

\emph{Case $N\equiv1\pmod8$.} Take $N=1$. Even $A$ give $A^2-1\equiv3\pmod4$. For $A=2t+1$, $A^2-1=4t(t+1)$, and by Lemma~\ref{lem:halving} $A$ passes iff $t(t+1)$ is a square modulo $2^{k_0}$, $k_0=m-2\ge1$. This condition depends only on $t\bmod 2^{k_0}$, while $t$ runs over $\Z/2^{m-1}\Z$; hence $C_m(1)=2D_{k_0}$ with
\[
D_{k_0}=\#\{t\in\Z/2^{k_0}\Z : t(t+1) \text{ is a square mod } 2^{k_0}\}.
\]
The involution $t\mapsto-1-t$ of $\Z/2^{k_0}\Z$ preserves $t(t+1)$ and exchanges the even and the odd residues, so $D_{k_0}=2E_{k_0}$, where $E_{k_0}$ counts the even $t$; thus $C_m(1)=4E_{m-2}$.

\emph{Claim:} $E_1=1$ and $E_{k_0}=s(k_0-2)$ for $k_0\ge2$. For $k_0=1$ the only even residue is $t=0$, which qualifies. Let $k_0\ge2$. The residue $t\equiv0$ qualifies. Every other even $t$ is $t=2^ju$ with $1\le j\le k_0-1$ and $u$ odd modulo $2^{k_0-j}$, and there are $2^{k_0-j-1}$ such $t$ for each $j$. Then $t(t+1)=2^jw$ with $w=u(1+2^ju)$ odd. By Lemma~\ref{lem:squares}, $t(t+1)$ is a square modulo $2^{k_0}$ iff $j$ is even and $w\equiv1\pmod{2^{\sigma}}$, $\sigma=\min(3,k_0-j)$. For even $j\ge2$ the map $u\mapsto u(1+2^ju)$ permutes the odd residues modulo $2^\sigma$: modulo $8$ it is $u\mapsto u+4$ for $j=2$ (as $4u^2\equiv4$) and the identity for $j\ge4$, and modulo $2^\sigma$ with $\sigma\le2$ it is the identity. Hence for each even $j$ exactly $2^{k_0-j-1}/2^{\sigma-1}=f(k_0-j)$ values of $t$ qualify, and
\[
E_{k_0}=1+\sum_{\substack{2\le j\le k_0-1\\ j\text{ even}}} f(k_0-j)=1+\sum_{0\le i<(k_0-2)/2}f(k_0-2-2i)=s(k_0-2),
\]
with $j=2i+2$ and Lemma~\ref{lem:sqcount}; for $k_0=2$ the sum is empty and $E_2=1=s(0)$. This proves the claim.

Therefore $C_3(1)=4E_1=4$ and $C_m(1)=4s(m-4)$ for $m\ge4$. For $m=4$ this is $4s(0)=4$; for $m\ge5$, Lemma~\ref{lem:sqcount} gives $4(2^{m-4}+8)/6=(2^{m-3}+16)/3$ for even $m$ and $4(2^{m-4}+10)/6=(2^{m-3}+20)/3$ for odd $m$. As a check at the smallest cases, $E_3=s(1)=2$, $E_4=s(2)=2$, $E_5=s(3)=3$, $E_6=s(4)=4$, $E_7=s(5)=7$, which agree with direct enumeration (Section~\ref{sec:numerics}); e.g.\ for $k_0=4$ the even $t\bmod 16$ with $t(t+1)$ a square are $t=0,4$, and for $k_0=5$ they are $t=0,16,20\bmod 32$. \qed

\begin{proof}[Proof of Corollary~\ref{cor:density}]
For $N\equiv3\pmod4$ the proportions are $1,\tfrac12,\tfrac14,\tfrac14,\dots$; for $N\equiv5\pmod 8$ they are $1,\tfrac12,\tfrac12,\tfrac14,\tfrac18,\tfrac18,\dots$. For $N\equiv1\pmod8$ they are $1,\tfrac12,\tfrac12,\tfrac14$ for $m\le4$, and $\pi_m=\tfrac1{24}+\tfrac{16}{3\cdot2^m}$ (even $m\ge6$) or $\tfrac1{24}+\tfrac{20}{3\cdot2^m}$ (odd $m\ge5$). All values exceed $1/24$, and $\pi_m\to1/24$ for $N\equiv1\pmod 8$.
\end{proof}

\begin{remark}
The limits are Haar measures: $\lim_m\pi_m(N)$ is the measure of $\{A\in\Z_2 : A^2-N\in\Z_2^2\}$, since passing at every level is equivalent to $A^2-N$ being a $2$-adic square (the set of squares of $\Z_2$ is closed). For $N=1$ and $A=2t+1$, the even-$t$ part has measure $\sum_{j\ge2\text{ even}}2^{-(j+1)}\cdot\tfrac14=\tfrac1{24}$ in the $t$-variable, the odd-$t$ part the same, and $A$ is odd with measure $\tfrac12$, which gives $\tfrac12\cdot\tfrac1{12}=\tfrac1{24}$.
\end{remark}

\section{The hyperbolic sieve modulo powers of two}

\begin{proof}[Proof of Theorem~\ref{thm:hyp}]
Since $k$ is odd it is a unit modulo $2^e$, and the substitution $x\mapsto kx$ in the definition gives $\Lset_{N,2^e,k}=\Lset_{kN,2^e,1}$. Put $K=kN$; it is odd. For $e=1$ both sides of the claim equal $1$ (the set is $\{0\}$ and $C_1=2$), so let $e\ge2$. Every element of $\Lset_{K,2^e,1}$ is even, being a sum of two units. For $a\in\Z$ the class of $a^2 \bmod 2^e$ depends only on $a\bmod 2^{e-1}$, because $(a+2^{e-1})^2=a^2+2^ea+2^{2e-2}$ and $2e-2\ge e$. We claim that for $a\in\Z$
\[
2a \bmod 2^e\in\Lset_{K,2^e,1}\iff a^2-K \text{ is a square modulo } 2^e .
\]
If $2a\equiv x+y$ with $xy\equiv K$, then $(x-a)^2=x^2-2ax+a^2\equiv x^2-x(x+y)+a^2=a^2-K\pmod{2^e}$. Conversely, if $a^2-K\equiv z^2$, put $x=a+z$, $y=a-z$; then $xy\equiv K$, so $x,y$ are units because $K$ is odd, and $x+y=2a$. Since $a\mapsto 2a$ maps $\Z/2^{e-1}\Z$ bijectively onto the even residues modulo $2^e$, $|\Lset_{K,2^e,1}|$ is the number of $a\bmod2^{e-1}$ with $a^2-K$ a square modulo $2^e$; by the periodicity just noted this is $C_e(K)/2$.

The listed values and the recursion follow from Theorem~\ref{thm:main} applied to $K=kN$; note that $kN\equiv3\pmod4$, $kN\equiv5\pmod8$ and $kN\equiv1\pmod8$ are equivalent to $N\equiv3k\pmod 4$, $N\equiv5k\pmod8$ and $N\equiv k\pmod8$ respectively, because $k^2\equiv1\pmod8$. For $N\equiv k\pmod 8$ and $e\ge5$, Theorem~\ref{thm:main} gives $2|\Lset_{N,2^e,k}|-|\Lset_{N,2^{e+1},k}|=2$ for even $e$ and $4$ for odd $e$, that is, $|\Lset_{N,2^{e+1},k}|=2(|\Lset_{N,2^e,k}|-2^{e\bmod2})$.
\end{proof}

\begin{remark}
Two different sets must not be confused. For odd $M$, \cite[Lemma~4.2]{hittmeir2023} identifies $\Lset_{N,M,k}$ with the full set of $x\bmod M$ for which $x^2-4kN$ is a square. Modulo $2^e$ the relevant comparison is with the half-centre $A=x/2$ and the polynomial $A^2-kN$: the pass set of $A^2-kN$ modulo $2^e$ has $C_e(kN)$ elements, and the hyperbolic set has half as many, because $x=2A$ determines $A$ only modulo $2^{e-1}$.
\end{remark}

\section{Sieve primes dividing the multiplier}

Algorithms that sieve $A^2-kn$ with several multipliers $k$ usually discard the primes that divide $kn$~\cite[Algorithm~4.3]{hittmeir2023}. A prime dividing $k$ exactly once still carries some information:

\begin{proposition}\label{prop:divk}
Let $\ell$ be an odd prime, $K\in\Z$ with $v_\ell(K)=1$, and $e\ge2$. Then the number of $A\bmod\ell^e$ for which $A^2-K$ is a square modulo $\ell^e$ is $\ell^e-\ell^{e-1}$.
\end{proposition}

\begin{proof}
If $\ell\nmid A$ then $A^2-K\equiv A^2\pmod\ell$ is a non-zero quadratic residue, hence a square modulo $\ell^e$ by Hensel's lemma. If $\ell\mid A$ then $v_\ell(A^2-K)=1<e$, and an element of valuation $1$ is not a square modulo $\ell^e$ for $e\ge2$, by the same valuation argument as in Lemma~\ref{lem:squares}.
\end{proof}

Thus such a prime removes a proportion $1/\ell$ at modulus $\ell^2$ and nothing more at higher powers, which is weaker than the factor of about $1/2$ obtained from a prime not dividing $kn$. For $\ell^2\mid K$ we do not give a general formula; an example is displayed in Section~\ref{sec:numerics}.

\section{Finite-range tests}\label{sec:numerics}
All stated identities were tested over the following finite ranges with the script \texttt{verify\_v11723.py} (pure Python~3; SHA-256 \texttt{64468db6\ldots aa9b95f}; exit status~$0$ iff all checks pass), pinned at \url{https://github.com/arsmagic/dyadic_shells/blob/c7a78b95f913a04f000ce7cd1593c6b652752049/research/twoadic-exact/verify_v11723.py}. These tests are not proofs. Results (all passed; full log \texttt{test\_log\_v11723.txt} in the same directory):
\begin{itemize}
\item Theorem~\ref{thm:main}, Corollary~\ref{cor:density} ($24C_m(N)>2^m$) and Proposition~\ref{prop:mod8}: odd $N\in[-259,259]$, $1\le m\le14$; $3640/3640$ cases each.
\item $E_{k_0}$ by direct enumeration for $k_0=1,\dots,7$: $1,1,2,2,3,4,7$; $E_{k_0}=s(k_0-2)$ for $k_0=2,\dots,7$; $C_m(1)=4E_{m-2}$ for $m=3,\dots,14$; Lemma~\ref{lem:sqcount} for $n=0,\dots,16$.
\item Theorem~\ref{thm:hyp}, $|\Lset_{N,2^e,k}|=C_e(kN)/2$: odd $N\in[-129,129]$, $k\in\{1,3,\dots,15\}$, $1\le e\le11$; $11440/11440$ triples.
\item The recursion of Theorem~\ref{thm:hyp} for $N\equiv k\pmod 8$, $5\le e\le12$, nine pairs $(N,k)$: $72/72$ steps with exponent $e\bmod2$, $0/72$ with $(e+1)\bmod2$. For instance $|\Lset_{1,2^e,1}|=4,4,6,8$ for $e=5,6,7,8$.
\item Example~4.7 of~\cite{hittmeir2023}: the values $2,4,4,6,8,14,24$ and $4,16,106$ are reproduced.
\item Odd prime powers, recursion of~\cite[Theorem~3.5]{hittmeir2023} with $(k^{-1}N\mid r)=1$, $r\le13$, $r^{e+1}\le30{,}000$: $983/983$ steps with exponent $e\bmod2$, $0/983$ with $(e+1)\bmod2$.
\item Proposition~\ref{prop:divk}: $\ell\in\{3,5,7\}$, $K=\ell n$ with $1\le n<30$, $\ell\nmid n$, $2\le e$, $\ell^e\le 200{,}000$; $345/345$ cases.
\item Example for $\ell^2\mid K$: $\ell=3$, $K=9$, counts $3,9,24,60,174,510$ modulo $3^e$, $e=1,\dots,6$ (proportions $1,1,0.889,0.741,0.716,0.700$).
\item Ingredients of the proofs: the bijection $u\mapsto u(1+2^ju)$ of odd residues modulo $2^\sigma$ for even $j\in[2,10]$, $\sigma\le3$ ($15/15$); Lemma~\ref{lem:halving} for $2\le n\le14$ and all $x\bmod 2^n$ ($32{,}764/32{,}764$); Lemma~\ref{lem:squares} against direct enumeration for $1\le n\le14$ ($32{,}766/32{,}766$).
\end{itemize}

\section{Formal verification}\label{sec:lean}
The following are proved in Lean~4 (v4.24.0) with Mathlib (tag \texttt{v4.24.0}) in the repository \url{https://github.com/arsmagic/dyadic_shells}, file \texttt{formal/DyadicShells/TwoAdicExact.lean}, commit \texttt{96b857a} on \texttt{main}; the continuous-integration build and axiom audit of that commit succeeded and report only the axioms \texttt{propext}, \texttt{Classical.choice}, \texttt{Quot.sound}:
\begin{itemize}
\item \texttt{DyadicShells.passCard\_unit\_sq\_mul} (Lemma~\ref{lem:scaling} for $R=\Z/M\Z$);
\item \texttt{DyadicShells.passCard\_eq\_of\_modEq\_eight} (Proposition~\ref{prop:mod8}).
\end{itemize}
They rely on the in-project lemma \texttt{DyadicShells.isSquare\_mod\_two\_pow\_of\_one\_mod\_eight} (an integer $c\equiv1\pmod 8$ is a square modulo every $2^k$), file \texttt{formal/DyadicShells/TwoAdicSaturation.lean}. A proposed contribution of the odd-unit criterion to Mathlib (\texttt{ZMod.isSquare\_two\_pow\_iff\_of\_odd} and two case forms) is under review in pull request~\#45540~\cite{mathlibpr} and is not merged. The closed forms of Theorem~\ref{thm:main}, the bound and limits of Corollary~\ref{cor:density}, Lemma~\ref{lem:sqcount}, Theorem~\ref{thm:hyp} and Proposition~\ref{prop:divk} are \emph{not} formally verified. Source: \url{https://github.com/arsmagic/dyadic_shells/blob/96b857add4a649220888d6e32f862dd96f255daa/formal/DyadicShells/TwoAdicExact.lean}.

\section*{Use of AI tools}
% Adapt to the target journal's policy before submission.
AI tools were used in deriving and checking the computations, in drafting the text and the test script, and in writing the Lean formalisation. The author is responsible for the content.

\begin{thebibliography}{99}
\bibitem{harveyhittmeir} D.~Harvey and M.~Hittmeir, \emph{A log-log speedup for exponent one-fifth deterministic integer factorisation}, Math. Comp. \textbf{91} (2022), no.~335, 1367--1379.
\bibitem{hittmeir2018} M.~Hittmeir, \emph{A babystep-giantstep method for faster deterministic integer factorization}, Math. Comp. \textbf{87} (2018), no.~314, 2915--2935.
\bibitem{hittmeir2023} M.~Hittmeir, \emph{Integer factorization as subset-sum problem}, J. Number Theory \textbf{249} (2023), 93--118, \href{https://doi.org/10.1016/j.jnt.2023.02.010}{doi:10.1016/j.jnt.2023.02.010}; cited from arXiv:2205.10074v1.
\bibitem{khach11517} A.~A. Khachatryan, \emph{Anchored sparse-offset screening of semiprimes}, Version 11.5.17, Zenodo, 2026, \href{https://doi.org/10.5281/zenodo.23432486}{doi:10.5281/zenodo.23432486}.
\bibitem{lawrence} F.~W. Lawrence, \emph{Factorisation of numbers}, Messenger of Math. \textbf{24} (1895), 100--109.
\bibitem{mathlibpr} A.~A. Khachatryan, \emph{Odd-unit square criterion in \texttt{ZMod (2\textasciicircum n)}}, pull request \#45540 to \texttt{leanprover-community/mathlib4}, open as of 2026-11-08, \url{https://github.com/leanprover-community/mathlib4/pull/45540}.
\bibitem{mckee} J.~McKee, \emph{Speeding Fermat's factoring method}, Math. Comp. \textbf{68} (1999), no.~228, 1729--1737.
\bibitem{serre} J.-P. Serre, \emph{A Course in Arithmetic}, Graduate Texts in Mathematics~7, Springer, New York, 1973.
\bibitem{stangl} W.~D. Stangl, \emph{Counting squares in $\mathbb{Z}_n$}, Math. Mag. \textbf{69} (1996), no.~4, 285--289.
\bibitem{stanica2026} P.~Stănică, \emph{Improved deterministic integer factorization by quadratic-image sieving}, arXiv:2611.06580v1, 2026.
\bibitem{shparlinski} I.~E. Shparlinski, \emph{Modular hyperbolas}, Jpn. J. Math. \textbf{7} (2012), 235--294; arXiv:1103.2879v4.
\bibitem{smh2026} P.~Stănică, E.~Mulder and M.~Hittmeir, \emph{Multi-target hyperbolic sieves and elliptic trace obstructions}, arXiv:2606.13018v2, 2026.
\end{thebibliography}

\end{document}
